% Section 12
\section{Appendix D: Consistency of the Double-Null Gauge Condition}
\label{gauge_consistency}

We verify that the double-null gauge condition $A_+ = A_- = 0$ is consistent with 
Maxwell's equations and leaves exactly two physical degrees of freedom.

\subsection{Coordinates and Notation}

In double-null coordinates $(x^+, x^-, x^1, x^2)$ the Minkowski metric reads
\begin{equation}
    ds^2 = -2\,dx^+\,dx^- + (dx^1)^2 + (dx^2)^2,
\end{equation}
so that $\partial_+ = \partial/\partial x^+$, $\partial_- = \partial/\partial x^-$, and 
the d'Alembertian is
\begin{equation}
    \Box = -2\partial_+\partial_- + \partial_i\partial^i, \qquad i = 1,2.
\end{equation}
The gauge field has components $A_\mu = (A_+, A_-, A_1, A_2)$ and transforms as
\begin{equation}
    A_\mu \to A_\mu + \partial_\mu\Lambda
\end{equation}
under a gauge transformation with parameter $\Lambda(x^+, x^-, x^1, x^2)$.

\subsection{Imposing the Gauge Conditions in Stages}

The single gauge parameter $\Lambda$ is a function of all four coordinates. We use 
it in three stages, each consuming a different functional dependence.

\paragraph{Stage 1: $A_+ = 0$.}
Requiring $A_+ + \partial_+\Lambda = 0$ determines $\Lambda$ as a function of $x^+$:
\begin{equation}
    \Lambda(x^+, x^-, x^i) = -\int^{x^+} A_+(s, x^-, x^i)\,ds 
                              + \Lambda_0(x^-, x^i),
\end{equation}
where $\Lambda_0(x^-, x^i)$ is an arbitrary residual function independent of $x^+$.

\paragraph{Stage 2: $A_- = 0$.}
Under the residual freedom $\Lambda_0(x^-, x^i)$, the $A_-$ component transforms as
$A_- \to A_- + \partial_-\Lambda_0$. Setting this to zero determines $\Lambda_0$ as a 
function of $x^-$:
\begin{equation}
    \Lambda_0(x^-, x^i) = -\int^{x^-} A_-(s, x^i)\,ds + \Lambda_1(x^i),
\end{equation}
where $\Lambda_1(x^i)$ is a residual function of the transverse coordinates only.

\paragraph{Stage 3: Transverse Coulomb condition $\partial^i A_i = 0$.}
Under the remaining freedom $\Lambda_1(x^i)$, the transverse components transform as 
$A_i \to A_i + \partial_i\Lambda_1$. Imposing $\partial^i A_i = 0$ requires
\begin{equation}
    \nabla^2_\perp\Lambda_1 = -\partial^i A_i,
\end{equation}
which is a Poisson equation in the transverse plane with a unique solution (up to 
harmonic functions, fixed by boundary conditions). This exhausts the gauge freedom 
entirely.

The three conditions have consumed the full gauge parameter $\Lambda$ in stages, 
with no over-determination: each stage uses a strictly different functional 
dependence of $\Lambda$.

\subsection{Physical Degrees of Freedom}

After imposing $A_+ = A_- = 0$ and $\partial^i A_i = 0$, the remaining components 
are $A_1$ and $A_2$ subject to one differential constraint $\partial^i A_i = 0$. 
In momentum space this reads $k^i A_i = 0$.

For our specific construction, the photons propagate in the $\pm z$ direction, 
so their wave-vectors are purely null: $k^i_\perp = 0$ for $i = 1,2$. The 
transverse Coulomb condition $k^i A_i = 0$ is therefore trivially satisfied for 
any $A_1, A_2$, and both transverse components are independent. The two circular 
polarization vectors
\begin{equation}
    \epsilon^i_{(+1)} = \frac{1}{\sqrt{2}}\begin{pmatrix}1\\i\end{pmatrix}, \qquad
    \epsilon^i_{(-1)} = \frac{1}{\sqrt{2}}\begin{pmatrix}1\\-i\end{pmatrix}
\end{equation}
are the two independent physical degrees of freedom, corresponding to helicity 
$\lambda = \pm 1$. This is the correct count for a massless spin-1 field in four 
dimensions.

The trivial satisfaction of the transverse Coulomb condition 
relies on the photons propagating purely in the $\pm z$ direction 
with vanishing transverse momentum $k^i_\perp = 0$. For photons 
with nonzero transverse momentum the condition $k^i A_i = 0$ 
imposes a non-trivial constraint and the degree of freedom count 
requires revisiting.

\subsection{Consistency with Maxwell's Equations}

With $A_+ = A_- = 0$ and $\partial^i A_i = 0$, the non-zero field strength 
components are
\begin{equation}
    F_{+i} = \partial_+ A_i, \qquad F_{-i} = \partial_- A_i, \qquad 
    F_{ij} = \partial_i A_j - \partial_j A_i, \qquad F_{+-} = 0.
\end{equation}
Maxwell's equations $\partial^\mu F_{\mu\nu} = J_\nu$ decompose as follows.

\paragraph{Transverse components $\nu = i$:}
\begin{equation}
    -2\partial_+\partial_- A_i + \partial^j(\partial_j A_i - \partial_i A_j) = J_i.
\end{equation}
Using $\partial^i A_i = 0$, the second term simplifies to $\partial^j\partial_j A_i$, 
giving the standard wave equation
\begin{equation}
    \boxed{\Box A_i = J_i.}
    \label{eq:wave}
\end{equation}

\paragraph{Null components $\nu = \pm$:}
\begin{equation}
    -\partial^i\partial_\pm A_i = J_\pm.
\end{equation}
These are constraint equations relating the null source components to the transverse 
field. For our bilocal source, which is purely transverse, $J_+ = J_- = 0$, and the 
constraints reduce to
\begin{equation}
    \partial^i\partial_\pm A_i = 0.
\end{equation}
Differentiating the transverse Coulomb condition $\partial^i A_i = 0$ with respect 
to $x^\pm$ gives $\partial^i(\partial_\pm A_i) = 0$, so these constraints are 
automatically satisfied. \checkmark

\subsection{Summary}

The double-null gauge condition $A_+ = A_- = 0$ supplemented by the transverse 
Coulomb condition $\partial^i A_i = 0$:
\begin{enumerate}
    \item uses the gauge freedom of $\Lambda$ completely and without over-determination,
          via a three-stage procedure consuming the $x^+$-, $x^-$-, and $x^i$-dependent
          parts of $\Lambda$ in sequence;
    \item leaves exactly two physical degrees of freedom, the helicity eigenstates
          $\lambda = \pm 1$, consistent with the known count for a massless spin-1 field;
    \item reduces Maxwell's equations to the standard transverse wave equation
          $\Box A_i = J_i$, with the null constraint equations automatically satisfied
          for a purely transverse source.
\end{enumerate}
The gauge is therefore internally consistent and does not over-constrain the physical 
degrees of freedom. The PCT wedge-swap symmetry is preserved throughout: 
conditions on $A_+$ and $A_-$ are imposed symmetrically, and neither null direction 
is privileged over the other.
